\documentclass[12pt]{article}
\usepackage[T1]{fontenc}
\usepackage[utf8]{inputenc}
\IfFileExists{newtxtext.sty}{
  \usepackage{newtxtext}
}{
  \IfFileExists{lmodern.sty}{\usepackage{lmodern}}{}
}
\usepackage[margin=1in]{geometry}
\usepackage{setspace}
\usepackage[hidelinks]{hyperref}

\onehalfspacing
\setlength{\parindent}{0pt}
\setlength{\parskip}{0.6em}

\begin{document}

\begin{center}
{\Large Public Co-Creation Hubs, Cross-Regional Pipelines, and Dynamic Efficiency in Thin Regional Ecosystems\par}
\vspace{0.5em}
{\large \"{O}ffentliche Co-Creation-Hubs, \"uberregionale Pipelines und dynamische Effizienz in d\"unnen regionalen \"Okosystemen\par}
\vspace{1em}
{\normalsize Original Paper\par}
\end{center}

\textbf{Author}\\
Ryota Matsuki

\textbf{Affiliation}\\
Independent Researcher, Matsuyama, Ehime, Japan

\textbf{Corresponding author}\\
Ryota Matsuki\\
E-mail: \href{mailto:ryota.matsuki@gmail.com}{ryota.matsuki@gmail.com}

\textbf{Keywords}\\
public co-creation hubs; cross-regional pipelines; regional innovation; place-based policy; thin regional ecosystems; dynamic efficiency

\textbf{JEL codes}\\
H70, O31, O38, R11, R12, R58

\textbf{Abstract}\\
Local governments increasingly support co-creation hubs as place-based instruments for strengthening regional innovation. Yet the welfare case for hub policy is unclear in thin regional ecosystems, where the local participant pool is limited and repeated local interaction may quickly become redundant. This paper develops a two-region, two-period framework in which a hub raises local matching in a target region, while pipeline policy creates cross-regional renewal through outside linkages. The model focuses on a wedge that arises when current hub services and directly brokered outside links are contractible, but the future novelty created inside the regional system is not. The framework delivers three conditional implications. First, decentralized support and planner evaluation need not coincide because renewal effort is underprovided and hub support can be either too weak or too strong. Second, a hub is socially efficient as a stand-alone intervention only under limited conditions. Third, pipeline policy has positive stand-alone value, but a hub-plus-pipeline bundle dominates pipeline policy alone only under specific parameter conditions. The framework generates empirically testable predictions on repeated ties, outside collaboration, and absorptive conditions, and it implies that hubs in thin regional ecosystems should be evaluated not only against no-policy baselines, but also against pipeline-only renewal strategies.

\textbf{Zusammenfassung}\\
Lokale Regierungen f\"ordern zunehmend Co-Creation-Hubs als ortsbezogene Instrumente zur St\"arkung regionaler Innovation. Ob Hub-Politik in d\"unnen regionalen \"Oko\-systemen wohlfahrts\-steigernd ist, bleibt jedoch unklar, weil der lokale Teilnehmerpool begrenzt ist und sich wiederholende lokale Interaktionen rasch redundant werden k\"onnen. Dieser Beitrag entwickelt ein Zwei-Regionen-Zwei-Perioden-Modell, in dem ein Hub das lokale Matching in einer Zielregion verbessert, w\"ahrend eine Pipeline-Politik durch externe Verbindungen \"uberregionale Erneuerung schafft. Im Mittelpunkt steht ein Keil, der entsteht, wenn laufende Hub-Leistungen und direkt vermittelte externe Verbindungen vertraglich erfassbar sind, die k\"unftig im regionalen System entstehende Neuheit jedoch nicht. Das Modell liefert drei bedingte Ergebnisse. Erstens k\"onnen dezentrale Unterst\"utzung und planerische Bewertung auseinanderfallen, weil Erneuerungsanstrengungen unterbereitgestellt werden und Hub-Unterst\"utzung sowohl zu schwach als auch zu stark ausfallen kann. Zweitens ist ein Hub als isolierte Ma{\ss}nahme nur unter begrenzten Bedingungen sozial effizient. Drittens hat Pipeline-Politik auch f\"ur sich genommen positiven Wert, doch ein B\"undel aus Hub und Pipeline dominiert Pipeline-Politik allein nur in bestimmten Parameterbereichen. Daraus ergeben sich empirisch pr\"ufbare Vorhersagen zu wiederholten Beziehungen, externer Zusammenarbeit und absorptiven Voraussetzungen. F\"ur d\"unne regionale \"Okosysteme folgt daraus, dass Hubs nicht nur gegen eine Situation ohne Politik, sondern auch gegen reine Pipeline-Strategien bewertet werden sollten.

\textbf{Acknowledgements}\\
No non-author acknowledgements are declared.

\textbf{Funding}\\
No specific funding was received for this research.

\textbf{Competing interests}\\
The author has no competing interests to declare.

\textbf{Author contributions}\\
The sole author conceived the study, developed the model, carried out the analysis, and wrote the manuscript.

\textbf{Code availability}\\
The symbolic verification scripts used for the manuscript revision are available from the author upon reasonable request.

\textbf{Use of AI tools}\\
During preparation of the manuscript, the author used ChatGPT (OpenAI) to improve wording and exposition in selected passages. The author independently verified all arguments and results and takes full responsibility for the final text.

\end{document}
